\documentclass[11pt,a4paper]{article}
\usepackage[margin=22mm]{geometry}
\usepackage{amsmath,amssymb,graphicx,booktabs,array}
\usepackage[colorlinks=true,linkcolor=black,urlcolor=blue]{hyperref}
\usepackage{fancyhdr}
\pagestyle{fancy}\fancyhf{}\fancyhead[L]{DDPG, TD3, and SAC: motivation and mathematics}\fancyfoot[C]{\thepage}
\setlength{\headheight}{14pt}
\setlength{\parindent}{0pt}\setlength{\parskip}{7pt}
\newcommand{\underE}[2]{\underset{\substack{#1}}{\mathrm E}\left[#2\right]}
\newsavebox{\eqbox}
\newcommand{\fitmath}[1]{\begin{equation}\sbox{\eqbox}{$\displaystyle #1$}\ifdim\wd\eqbox>.94\linewidth\resizebox{.94\linewidth}{!}{\usebox{\eqbox}}\else\usebox{\eqbox}\fi\end{equation}}
\newcommand{\motivation}[1]{\subsection{Why this algorithm?}\noindent\fbox{\begin{minipage}{.94\linewidth}#1\end{minipage}}\par}
\begin{document}
\begin{center}{\LARGE\bfseries DDPG, TD3, and SAC}\\[6pt]{\large Motivation and mathematical foundations}\end{center}
\textbf{Reading guide.} Each algorithm starts with its motivation, followed by equations in source order with short explanations. This is a mathematical study note: implementation code, pseudocode instructions, isolated symbol lists, and repeated API-documentation copies are omitted. Repeated targets in the update rules are referenced rather than printed again. Equation numbering is added for navigation; mathematical notation follows Spinning Up.

\textbf{Sources.} OpenAI Spinning Up documentation, retrieved 18 October 2026. Each algorithm section links to its source page. Motivation summaries paraphrase those pages.
\subsection*{Notation used throughout}
\begin{tabular}{@{}p{.25\linewidth}p{.69\linewidth}@{}}\toprule
Symbol & Meaning\\\midrule
$s,a,r,s'$ & State, action, observed reward, and next state.\\
$\gamma,d$ & Discount factor and terminal indicator; $d=1$ means the next state is terminal.\\
$\mathcal D,B$ & Replay buffer and sampled minibatch; $|B|$ is the minibatch size.\\
$\theta,\phi_i$ & Policy parameters and critic parameters. A single critic uses $\phi$.\\
$\mathrm{targ},\rho$ & Target-network subscript and Polyak averaging coefficient, usually close to 1.\\
$\mu_\theta,\pi_\theta$ & Deterministic policy and stochastic policy, respectively.\\
$\alpha,H$ & Entropy coefficient and entropy.\\\bottomrule
\end{tabular}
\tableofcontents


\clearpage

\section{Deep Deterministic Policy Gradient}

\textbf{Source:} \url{https://spinningup.openai.com/en/latest/algorithms/ddpg.html}

\motivation{\textbf{Problem: continuous-action maximization.} Q-learning needs the action that maximizes a learned value function. Enumerating actions is practical for a small discrete action set, but impossible for a continuous one; repeatedly solving a numerical optimization problem is expensive.

\textbf{DDPG's idea.} Learn a deterministic actor that approximates the maximizing action, using gradients through the critic. Learn the critic from off-policy transitions. Replay reuses experience, while slowly changing target networks reduce instability in the Bellman targets.}

\subsection{Equations and their roles}

\subsubsection*{Choosing an optimal action}

The optimal action maximizes the optimal action-value function.

\fitmath{a^*(s) = \arg \max_a Q^*(s,a).}

\subsubsection*{Replacing action optimization with an actor}

The actor approximates the maximizer, avoiding a separate optimization at each decision.

\fitmath{\max_a Q(s,a) \approx Q(s,\mu(s))}

\subsubsection*{Optimal Bellman equation}

Here $s'\sim P$ abbreviates sampling from the transition distribution $P(\cdot\mid s,a)$.

\fitmath{Q^*(s,a) = \underset{s' \sim P}{{\mathrm E}}\left[r(s,a) + \gamma \max_{a'} Q^*(s', a')\right]}

\subsubsection*{Initial mean-squared Bellman error}

This preliminary loss uses the same critic in the prediction and target, which can make learning unstable.

\fitmath{L(\phi, {\mathcal D}) = \underset{(s,a,r,s',d) \sim {\mathcal D}}{{\mathrm E}}\left[
    \Bigg( Q_{\phi}(s,a) - \left(r + \gamma (1 - d) \max_{a'} Q_{\phi}(s',a') \right) \Bigg)^2
    \right]}

\subsubsection*{The bootstrap target}

The following is the target term inside that loss; it is an expression, not a separate optimization objective.

\fitmath{r + \gamma (1 - d) \max_{a'} Q_{\phi}(s',a')}

\subsubsection*{A slowly changing target critic}

Polyak averaging makes the target critic track the learned critic gradually.

\fitmath{\phi_{\text{targ}} \leftarrow \rho \phi_{\text{targ}} + (1 - \rho) \phi,}

\subsubsection*{DDPG critic loss}

The practical loss uses a target critic and a target actor. The terminal factor removes future value at terminal transitions.

\fitmath{L(\phi, {\mathcal D}) = \underset{(s,a,r,s',d) \sim {\mathcal D}}{{\mathrm E}}\left[
    \Bigg( Q_{\phi}(s,a) - \left(r + \gamma (1 - d) Q_{\phi_{\text{targ}}}(s', \mu_{\theta_{\text{targ}}}(s')) \right) \Bigg)^2
    \right],}

\subsubsection*{Actor objective}

Maximize the critic value of the actor action, treating critic parameters as fixed.

\fitmath{\max_{\theta} \underset{s \sim {\mathcal D}}{{\mathrm E}}\left[ Q_{\phi}(s, \mu_{\theta}(s)) \right].}

\subsubsection*{Exploration action}

Training adds Gaussian exploration noise and clips actions to the valid range. The source leaves the Gaussian parameters unspecified here.

\fitmath{a = \text{clip}(\mu_{\theta}(s) + \epsilon, a_{Low}, a_{High}),\qquad \epsilon\sim\mathcal N}

\subsubsection*{Minibatch target}

This target is used for the critic gradient below.

\fitmath{y(r,s',d) = r + \gamma (1-d) Q_{\phi_{\text{targ}}}(s', \mu_{\theta_{\text{targ}}}(s'))}

\subsubsection*{Critic gradient}

Take a descent step using this minibatch gradient.

\fitmath{\nabla_{\phi} \frac{1}{|B|}\sum_{(s,a,r,s',d) \in B} \left( Q_{\phi}(s,a) - y(r,s',d) \right)^2}

\subsubsection*{Actor gradient}

Take an ascent step using this minibatch gradient.

\fitmath{\nabla_{\theta} \frac{1}{|B|}\sum_{s \in B}Q_{\phi}(s, \mu_{\theta}(s))}

\subsubsection*{Target-network updates}

Initialize target networks from their corresponding learned networks, then apply these updates.

\fitmath{\begin{aligned}
                    \phi_{\text{targ}} &\leftarrow \rho \phi_{\text{targ}} + (1-\rho) \phi \\
                    \theta_{\text{targ}} &\leftarrow \rho \theta_{\text{targ}} + (1-\rho) \theta
                \end{aligned}}

\clearpage

\section{Twin Delayed DDPG}

\textbf{Source:} \url{https://spinningup.openai.com/en/latest/algorithms/td3.html}

\motivation{\textbf{Problem: the actor exploits critic errors.} DDPG can overestimate action values. Its actor can then select actions because their estimated values are wrong, producing brittle behavior.

\textbf{TD3's three changes.} Use the smaller of two target critics to limit overestimation; update the actor and target networks less frequently than the critics; and smooth target actions with clipped noise to discourage exploitation of narrow, erroneous peaks.}

\subsection{Equations and their roles}

\subsubsection*{Target policy smoothing}

Clip noise to $[-c,c]$, add it to the target action, and clip the result to the action bounds. The Gaussian notation is retained exactly as on the page.

\fitmath{a'(s') = \text{clip}\left(\mu_{\theta_{\text{targ}}}(s') + \text{clip}(\epsilon,-c,c), a_{Low}, a_{High}\right), \;\;\;\;\; \epsilon \sim \mathcal{N}(0, \sigma)}

\subsubsection*{Clipped double-Q target}

Use the smaller target-critic value as the shared backup.

\fitmath{y(r,s',d) = r + \gamma (1 - d) \min_{i=1,2} Q_{\phi_{i, \text{targ}}}(s', a'(s')),}

\subsubsection*{First critic loss}

Both critics regress toward the same target.

\fitmath{L(\phi_1, {\mathcal D}) = \underE{(s,a,r,s',d) \sim {\mathcal D}}{
    \Bigg( Q_{\phi_1}(s,a) - y(r,s',d) \Bigg)^2
    },}

\subsubsection*{Second critic loss}

The second critic has separate parameters but uses the same backup.

\fitmath{L(\phi_2, {\mathcal D}) = \underE{(s,a,r,s',d) \sim {\mathcal D}}{
    \Bigg( Q_{\phi_2}(s,a) - y(r,s',d) \Bigg)^2
    }.}

\subsubsection*{Actor objective}

The actor maximizes the first critic.

\fitmath{\max_{\theta} \underset{s \sim {\mathcal D}}{{\mathrm E}}\left[ Q_{\phi_1}(s, \mu_{\theta}(s)) \right],}

\subsubsection*{Exploration action}

As in DDPG, training actions include exploration noise. This is distinct from noise used to smooth target actions.

\fitmath{a = \text{clip}(\mu_{\theta}(s) + \epsilon, a_{Low}, a_{High}),\qquad \epsilon\sim\mathcal N}

\subsubsection*{Critic gradients}

Take descent steps for both critics using the target above.

\fitmath{\begin{aligned}
                    & \nabla_{\phi_i} \frac{1}{|B|}\sum_{(s,a,r,s',d) \in B} \left( Q_{\phi_i}(s,a) - y(r,s',d) \right)^2 && \text{for } i=1,2
                \end{aligned}}

\subsubsection*{Delayed actor gradient}

Take an ascent step less frequently than the critic updates. Spinning Up describes one policy update for every two critic updates as the paper recommendation.

\fitmath{\nabla_{\theta} \frac{1}{|B|}\sum_{s \in B}Q_{\phi_1}(s, \mu_{\theta}(s))}

\subsubsection*{Delayed target-network updates}

Initialize targets from their learned counterparts. Apply these averages only when the actor is updated.

\fitmath{\begin{aligned}
                        \phi_{\text{targ},i} &\leftarrow \rho \phi_{\text{targ}, i} + (1-\rho) \phi_i && \text{for } i=1,2\\
                        \theta_{\text{targ}} &\leftarrow \rho \theta_{\text{targ}} + (1-\rho) \theta
                    \end{aligned}}

\clearpage

\section{Soft Actor-Critic}

\textbf{Source:} \url{https://spinningup.openai.com/en/latest/algorithms/sac.html}

\motivation{\textbf{Problem: premature concentration of the policy.} Optimizing reward alone can concentrate a policy too early on poor behavior. An entropy bonus rewards randomness, encouraging exploration and helping avoid poor local optima.

\textbf{SAC's idea.} Learn a stochastic actor and off-policy critics using a reward-plus-entropy objective. SAC shares clipped double-Q learning with TD3, but was developed roughly concurrently with it: it is not a direct successor built to repair TD3. Stochastic actions provide an effect similar to target policy smoothing; no target actor is required in the version presented here.}

\subsection{Equations and their roles}

\subsubsection*{Entropy}

Entropy measures randomness under a probability mass or density function.

\fitmath{H(P) = \underE{x \sim P}{-\log P(x)}.}

\subsubsection*{Entropy-regularized objective}

The trajectory $\tau$ follows the policy. The positive coefficient $\alpha$ trades off reward and entropy.

\fitmath{\pi^* = \arg \max_{\pi} \underE{\tau \sim \pi}{ \sum_{t=0}^{\infty} \gamma^t \bigg( R(s_t, a_t, s_{t+1}) + \alpha H\left(\pi(\cdot|s_t)\right) \bigg)},}

\subsubsection*{State value with entropy}

The state value includes entropy bonuses from the first timestep onward.

\fitmath{V^{\pi}(s) = \underE{\tau \sim \pi}{ \left. \sum_{t=0}^{\infty} \gamma^t \bigg( R(s_t, a_t, s_{t+1}) + \alpha H\left(\pi(\cdot|s_t)\right) \bigg) \right| s_0 = s}}

\subsubsection*{Action value with entropy}

The action value excludes the entropy bonus at the first timestep. This convention explains the difference from the state value.

\fitmath{Q^{\pi}(s,a) = \underE{\tau \sim \pi}{ \left. \sum_{t=0}^{\infty} \gamma^t  R(s_t, a_t, s_{t+1}) + \alpha \sum_{t=1}^{\infty} \gamma^t H\left(\pi(\cdot|s_t)\right)\right| s_0 = s, a_0 = a}}

\subsubsection*{Connecting state and action values}

Average over the current action, then add the current-state entropy bonus.

\fitmath{V^{\pi}(s) = \underE{a \sim \pi}{Q^{\pi}(s,a)} + \alpha H\left(\pi(\cdot|s)\right)}

\subsubsection*{Entropy-regularized Bellman equation}

The next-state value includes both future return and next-state policy entropy.

\fitmath{\begin{aligned}Q^{\pi}(s,a) &= \underE{s' \sim P \\ a' \sim \pi}{R(s,a,s') + \gamma\left(Q^{\pi}(s',a') + \alpha H\left(\pi(\cdot|s')\right) \right)} \\
&= \underE{s' \sim P}{R(s,a,s') + \gamma V^{\pi}(s')}.\end{aligned}}

\subsubsection*{Writing entropy as a log-probability}

Replacing expected entropy with expected negative log-probability yields the form used in learning.

\fitmath{\begin{aligned}Q^{\pi}(s,a) &= \underE{s' \sim P \\ a' \sim \pi}{R(s,a,s') + \gamma\left(Q^{\pi}(s',a') + \alpha H\left(\pi(\cdot|s')\right) \right)} \\
&= \underE{s' \sim P \\ a' \sim \pi}{R(s,a,s') + \gamma\left(Q^{\pi}(s',a') - \alpha \log \pi(a'|s') \right)}\end{aligned}}

\subsubsection*{A sampled Bellman backup}

Sample $\tilde a'$ freshly from the policy; the reward and next state come from replay.

\fitmath{\begin{aligned}Q^{\pi}(s,a) &\approx r + \gamma\left(Q^{\pi}(s',\tilde{a}') - \alpha \log \pi(\tilde{a}'|s') \right), \;\;\;\;\;  \tilde{a}' \sim \pi(\cdot|s').\end{aligned}}

\subsubsection*{Critic loss}

Each critic regresses toward a shared target.

\fitmath{L(\phi_i, {\mathcal D}) = \underset{(s,a,r,s',d) \sim {\mathcal D}}{{\mathrm E}}\left[
    \Bigg( Q_{\phi_i}(s,a) - y(r,s',d) \Bigg)^2
    \right],}

\subsubsection*{SAC target}

The minimum of two target critics limits overestimation. The entropy contribution is $-\alpha\log\pi_\theta$; next actions come from the current policy.

\fitmath{y(r, s', d) = r + \gamma (1 - d) \left( \min_{j=1,2} Q_{\phi_{\text{targ},j}}(s', \tilde{a}') - \alpha \log \pi_{\theta}(\tilde{a}'|s') \right), \;\;\;\;\; \tilde{a}' \sim \pi_{\theta}(\cdot|s').}

\subsubsection*{Policy value objective}

The actor seeks both high value and high entropy.

\fitmath{\begin{aligned}V^{\pi}(s) &= \underE{a \sim \pi}{Q^{\pi}(s,a)} + \alpha H\left(\pi(\cdot|s)\right) \\
&= \underE{a \sim \pi}{Q^{\pi}(s,a) - \alpha \log \pi(a|s)}.\end{aligned}}

\subsubsection*{Reparameterized action}

The Gaussian mean and standard deviation are $\mu_\theta(s)$ and $\sigma_\theta(s)$. The operator $\odot$ is elementwise multiplication; $\tanh$ bounds the action.

\fitmath{\tilde{a}_{\theta}(s, \xi) = \tanh\left( \mu_{\theta}(s) + \sigma_{\theta}(s) \odot \xi \right), \;\;\;\;\; \xi \sim \mathcal{N}(0, I).}

\subsubsection*{Moving randomness outside the policy}

Reparameterization permits differentiation through the sampled action.

\fitmath{\underE{a \sim \pi_{\theta}}{Q^{\pi_{\theta}}(s,a) - \alpha \log \pi_{\theta}(a|s)} = \underE{\xi \sim \mathcal{N}}{Q^{\pi_{\theta}}(s,\tilde{a}_{\theta}(s,\xi)) - \alpha \log \pi_{\theta}(\tilde{a}_{\theta}(s,\xi)|s)}}

\subsubsection*{Practical actor objective}

Use the smaller learned critic in the entropy-regularized objective. This page uses a fixed entropy coefficient.

\fitmath{\begin{aligned}\max_{\theta} \underE{s \sim \mathcal{D} \\ \xi \sim \mathcal{N}}{\min_{j=1,2} Q_{\phi_j}(s,\tilde{a}_{\theta}(s,\xi)) - \alpha \log \pi_{\theta}(\tilde{a}_{\theta}(s,\xi)|s)},\end{aligned}}

\subsubsection*{Minibatch critic gradients}

Take descent steps for both critics using the SAC target above.

\fitmath{\begin{aligned}
                    & \nabla_{\phi_i} \frac{1}{|B|}\sum_{(s,a,r,s',d) \in B} \left( Q_{\phi_i}(s,a) - y(r,s',d) \right)^2 && \text{for } i=1,2
                \end{aligned}}

\subsubsection*{Minibatch actor gradient}

Take an ascent step. Here $\tilde a_\theta(s)$ denotes a differentiable action sample from the current policy, using reparameterization.

\fitmath{\nabla_{\theta} \frac{1}{|B|}\sum_{s \in B} \Big(\min_{i=1,2} Q_{\phi_i}(s, \tilde{a}_{\theta}(s)) - \alpha \log \pi_{\theta} \left(\left. \tilde{a}_{\theta}(s) \right| s\right) \Big),}

\subsubsection*{Target-critic updates}

Initialize target critics from learned critics, then apply Polyak averaging. SAC has no target actor in this formulation.

\fitmath{\begin{aligned}
                    \phi_{\text{targ},i} &\leftarrow \rho \phi_{\text{targ}, i} + (1-\rho) \phi_i && \text{for } i=1,2
                \end{aligned}}

\end{document}
